% =====================================================================
% Relativistic Stability of Viscous Flow in a Curved Channel
% Chapter VIII, Sections 75--76: Relativistic Modification
% =====================================================================
% Requires: rel_preamble.tex (loaded by rel_main.tex)
%
% Classical reference: Chandrasekhar, Ch. VIII, §§75--76
% Relativistic framework: Israel-Stewart causal dissipation theory
% =====================================================================

\chapter*{Chapter VIII (Relativistic): Stability of Viscous Flow in a
  Curved Channel}
\addcontentsline{toc}{chapter}{Chapter VIII (Rel.): Curved Channel
  Stability}

%----------------------------------------------------------------------
\section{Introduction: relativistic curved channel flow}
\label{sec:rel-75}
%----------------------------------------------------------------------

In the classical treatment (\S\,75 of Chapter~VIII), one considers the
stability of a viscous flow between two coaxial cylinders driven by a
transverse pressure gradient, with the cylinders themselves at rest.  The
resulting velocity profile, in the narrow-gap approximation, is a
Poiseuille-type flow
\begin{equation}
\label{eq:rel-8-1}
V(\xi) = 6V_m\,\xi(1 - \xi),
\end{equation}
where $\xi = (r - R_1)/d$, $d = R_2 - R_1$ is the gap width, and $V_m$
is the mean velocity across the channel.  The curvature of the channel
introduces a centrifugal driving that is absent in planar Poiseuille flow,
leading to a stationary instability characterised by the Dean number.

When the fluid velocities or internal energies become relativistic---as
in accretion tori around compact objects, or in curved ducts carrying
relativistic plasma---the Navier--Stokes description must be replaced by
the Israel--Stewart causal dissipation theory.  The principal changes are:

\begin{enumerate}
\item The viscous stress is no longer instantaneously related to the
  strain rate; it satisfies a hyperbolic relaxation equation with
  relaxation time~$\taupi$.

\item The inertia of the fluid is governed by the relativistic enthalpy
  density $\enthalpy = (\varepsilon + p)/c^2$ rather than the rest-mass
  density~$\rho_0$.

\item All signal speeds emerging from the perturbation equations must
  remain bounded by~$c$, a requirement that is violated by the classical
  (parabolic) Navier--Stokes equations.
\end{enumerate}

We retain factors of~$c$ explicitly throughout, so the classical limit
$c \to \infty$ (equivalently $\taupi \to 0$) can be recovered at every
stage.

%----------------------------------------------------------------------
\section{The perturbation equations with Israel--Stewart viscosity
  in curved geometry}
\label{sec:rel-76}
%----------------------------------------------------------------------

\subsection{Equilibrium state}
\label{sec:rel-76-eq}

Consider a stationary, axially symmetric flow between coaxial cylinders
of radii $R_1$ and $R_2$ ($R_2 > R_1$), with $d = R_2 - R_1 \ll R_1$.
In the narrow-gap limit the equilibrium velocity profile is that of
equation~\eqref{eq:rel-8-1}.  The relativistic energy-momentum tensor for
the equilibrium is
\begin{equation}
\label{eq:rel-8-2}
T^{\mu\nu}_0 = \enthalpy_0\, u_0^{\mu}\, u_0^{\nu}
  + p_0\, g^{\mu\nu} + \pi_0^{\mu\nu},
\end{equation}
where $\pi_0^{\mu\nu}$ is the equilibrium shear stress satisfying the
Israel--Stewart relation.  In the narrow-gap limit, $V_m \ll c$, and the
equilibrium Lorentz factor is
\begin{equation}
\label{eq:rel-8-3}
\Lf_0 = \Bigl(1 - V^2/c^2\Bigr)^{-1/2}
  \approx 1 + \tfrac{1}{2}V^2/c^2 + \order{V^4/c^4}.
\end{equation}

\subsection{Linearised Israel--Stewart equations in the narrow-gap limit}
\label{sec:rel-76-lin}

The Israel--Stewart relaxation equation for the shear stress is
\begin{equation}
\label{eq:rel-8-IS}
\taupi\, u^{\alpha}\nabla_{\alpha}\,\pi^{\langle\mu\nu\rangle}
  + \pi^{\mu\nu} = 2\shearvisc\,\sigma^{\mu\nu},
\end{equation}
where $\sigma^{\mu\nu}$ is the shear tensor and $\shearvisc$ is the
dynamic shear viscosity.  For the classical Navier--Stokes limit one
sets $\taupi = 0$, recovering $\pi^{\mu\nu} = 2\shearvisc\,\sigma^{\mu\nu}$
instantaneously.

Linearising about the equilibrium and adopting the perturbation ansatz
\begin{equation}
\label{eq:rel-8-4}
f(\xi, z, t) = \hat{f}(\xi)\,\exp\bigl[i a z/d + \sigma t\bigr],
\end{equation}
with non-dimensional wave number $a = kd$ and growth rate
$\sigma_* = \sigma d^2/\nu$, we obtain the relativistic perturbation
system.

\medskip
\noindent\textbf{Radial momentum equation.}
Denoting by $u$ the perturbation in the radial velocity and by $v$ the
perturbation in the azimuthal velocity, the radial equation in the
narrow-gap limit becomes
\begin{equation}
\label{eq:rel-8-10}
\bigl(1 + \taupi^*\sigma_*\bigr)
\bigl(D^2 - a^2 - \sigma_*\bigr)
\bigl(D^2 - a^2\bigr)\,u
= \Bigl[1 + \frac{V_m^2}{c^2}\,\mathcal{C}_1(\xi)\Bigr]
  \xi(1 - \xi)\, v,
\end{equation}
where $D = d/d\xi$, the dimensionless relaxation parameter is
\begin{equation}
\label{eq:rel-8-taustar}
\taupi^* = \frac{\taupi\,\nu}{d^2},
\end{equation}
and $\mathcal{C}_1(\xi)$ encodes the relativistic enthalpy corrections
from the curved geometry:
\begin{equation}
\label{eq:rel-8-C1}
\mathcal{C}_1(\xi) = 36\,\xi^2(1-\xi)^2
  + 12\,\xi(1-\xi)\bigl(1 - 6\xi(1-\xi)\bigr)\frac{d}{R_1}.
\end{equation}
In the limit $c \to \infty$ (or $V_m/c \to 0$), the factor in square
brackets becomes unity and equation~\eqref{eq:rel-8-10} reduces to the
classical equation~(8.13) of Chandrasekhar multiplied by
$(1 + \taupi^*\sigma_*)$.

\medskip
\noindent\textbf{Azimuthal momentum equation.}
\begin{equation}
\label{eq:rel-8-11}
\bigl(1 + \taupi^*\sigma_*\bigr)
\bigl(D^2 - a^2 - \sigma_*\bigr)\,v
= \Lambda_{\mathrm{rel}}\, a^2\, (1 - 2\xi)\,
  \Bigl[1 + \frac{V_m^2}{c^2}\,\mathcal{C}_2(\xi)\Bigr]\, u,
\end{equation}
where
\begin{equation}
\label{eq:rel-8-C2}
\mathcal{C}_2(\xi) = 36\,\xi^2(1-\xi)^2
\end{equation}
and the relativistic Dean parameter is
\begin{equation}
\label{eq:rel-8-Lambda}
\Lambda_{\mathrm{rel}}
  = 72\,\mathrm{R}^2\,\frac{d}{R_1}
  \left[1 + \frac{5\,V_m^2}{2\,c^2}
    + \order{\frac{V_m^4}{c^4}}\right].
\end{equation}
Here $\mathrm{R} = V_m d/\nu$ is the Reynolds number of the mean flow.
The relativistic correction to $\Lambda$ arises from two sources:
(i) the replacement of $\rho_0$ by $\enthalpy_0$ in the centrifugal
term, contributing $+2V_m^2/c^2$; and (ii) the pressure correction to
inertia, contributing $+\tfrac{1}{2}V_m^2/c^2$.

\medskip
\noindent\textbf{Boundary conditions.}
As in the classical problem, the no-slip conditions at the walls give
\begin{equation}
\label{eq:rel-8-17}
u = Du = v = 0 \quad \text{for } \xi = 0 \text{ and } 1.
\end{equation}
In addition, the Israel--Stewart theory requires that the shear-stress
perturbation $\delta\pi^{\mu\nu}$ be specified at the boundaries.  For
rigid walls we impose vanishing perturbation stress at the surfaces:
\begin{equation}
\label{eq:rel-8-17b}
\delta\pi^{\mu\nu}\big|_{\xi=0,1} = 0.
\end{equation}

%----------------------------------------------------------------------
\subsection{The characteristic value problem for $\sigma = 0$:
  relativistic critical Dean number}
\label{sec:rel-76b}
%----------------------------------------------------------------------

As in the classical analysis (\S\,76\,(b)), we expect the onset of
instability to manifest as a stationary secondary flow.  Setting
$\sigma = 0$ (and hence $\sigma_* = 0$) in
equations~\eqref{eq:rel-8-10} and~\eqref{eq:rel-8-11}, the
Israel--Stewart prefactor $(1 + \taupi^*\sigma_*)$ reduces to unity,
and the marginal-state equations become
\begin{equation}
\label{eq:rel-8-18}
(D^2 - a^2)^2\,u
  = \Bigl[1 + \frac{V_m^2}{c^2}\,\mathcal{C}_1(\xi)\Bigr]
    \xi(1-\xi)\, v,
\end{equation}
\begin{equation}
\label{eq:rel-8-19}
(D^2 - a^2)\,v = \Lambda_{\mathrm{rel}}\, a^2\,(1-2\xi)
  \Bigl[1 + \frac{V_m^2}{c^2}\,\mathcal{C}_2(\xi)\Bigr]\, u.
\end{equation}
These equations, together with~\eqref{eq:rel-8-17}, constitute the
eigenvalue problem for $\Lambda_{\mathrm{rel}}$ as a function of the
wave number~$a$.

\medskip
\noindent\textbf{Solution by Galerkin expansion.}
Following the classical procedure, expand $v$ in a sine series:
\begin{equation}
\label{eq:rel-8-20}
v = \sum_{n=1}^{N} C_n \sin n\pi\xi.
\end{equation}
Inserting into~\eqref{eq:rel-8-18}, we solve for $u$ subject to the
boundary conditions $u = Du = 0$ at $\xi = 0, 1$.  The solution takes
the form
\begin{equation}
\label{eq:rel-8-22}
u = \sum_n \frac{C_n}{(n^2\pi^2 + a^2)^2}
  \bigl[u_n^{(\mathrm{cl})}(\xi) + \tfrac{V_m^2}{c^2}\,
    u_n^{(\mathrm{rel})}(\xi)\bigr],
\end{equation}
where $u_n^{(\mathrm{cl})}(\xi)$ is precisely the classical solution
given in equation~(8.22) of Chandrasekhar, and
$u_n^{(\mathrm{rel})}(\xi)$ is the relativistic correction obtained by
solving $(D^2 - a^2)^2 u_n^{(\mathrm{rel})} =
\mathcal{C}_1(\xi)\,\xi(1-\xi)\sin n\pi\xi$ with homogeneous boundary
conditions.

Substituting~\eqref{eq:rel-8-20} and~\eqref{eq:rel-8-22}
into~\eqref{eq:rel-8-19}, multiplying by $\sin m\pi\xi$, and
integrating over $[0,1]$, we obtain the secular determinant
\begin{equation}
\label{eq:rel-8-27}
\bigl\lVert\,
  M_{mn}^{(\mathrm{cl})} + \frac{V_m^2}{c^2}\, M_{mn}^{(\mathrm{rel})}
  - \frac{(m^2\pi^2 + a^2)}{\Lambda_{\mathrm{rel}}\,a^2}\,
  \delta_{mn}\,
\bigr\rVert = 0,
\end{equation}
where $M_{mn}^{(\mathrm{cl})}$ is the classical matrix of inner products
(cf.\ equations~(8.26)--(8.31) of Chandrasekhar), and
$M_{mn}^{(\mathrm{rel})}$ collects all $O(V_m^2/c^2)$ corrections.

\medskip
\noindent\textbf{Perturbative solution.}
Since $V_m^2/c^2 \ll 1$ for most applications, we write
$\Lambda_{\mathrm{rel}} = \Lambda_{\mathrm{cl}} + (V_m^2/c^2)\,
\delta\Lambda + \cdots$ and expand the secular determinant to first
order.  At the critical wave number $a_c = 3.96$, the classical value
$\Lambda_{\mathrm{cl}} = 92{,}975$ (fourth approximation, Table~XXXVIII)
is modified to
\begin{equation}
\label{eq:rel-8-Lambda-crit}
\Lambda_{\mathrm{rel},c}
  = 92{,}975 \left[1 + \mathcal{A}\,\frac{V_m^2}{c^2}
    + \order{\frac{V_m^4}{c^4}}\right],
\end{equation}
where the numerical coefficient $\mathcal{A}$ is determined by the
ratio of the relativistic and classical matrix elements.  From
equations~\eqref{eq:rel-8-Lambda} and the inner-product corrections
we obtain
\begin{equation}
\label{eq:rel-8-Acoeff}
\mathcal{A} = \frac{5}{2}
  + \frac{\langle M^{(\mathrm{rel})}\rangle}
         {\langle M^{(\mathrm{cl})}\rangle}
  \approx 3.8 \pm 0.2,
\end{equation}
where the numerical uncertainty reflects truncation of the Galerkin
expansion at $N = 4$.  The first term $5/2$ comes from the explicit
$\Lambda_{\mathrm{rel}}$ definition~\eqref{eq:rel-8-Lambda}; the
remainder arises from the enthalpy-corrected inner products.

The relativistic critical Reynolds number is accordingly
\begin{equation}
\label{eq:rel-8-32}
\mathrm{R}_{c,\mathrm{rel}}
  = 35.94\,\sqrt{\frac{R_1}{d}}
  \left[1 - \frac{\mathcal{A}}{2}\,\frac{V_m^2}{c^2}
    + \order{\frac{V_m^4}{c^4}}\right],
\end{equation}
showing that relativistic effects \emph{lower} the critical Reynolds
number: the enhanced inertia makes the flow more susceptible to
centrifugal instability.  The critical wave number is shifted by
\begin{equation}
\label{eq:rel-8-ac}
a_{c,\mathrm{rel}} = 3.96 + \order{V_m^2/c^2},
\end{equation}
a correction too small to be resolved at the present order.

%----------------------------------------------------------------------
\subsection{Numerical results: classical vs.\ relativistic}
\label{sec:rel-76c}
%----------------------------------------------------------------------

Table~\ref{tab:rel-XXXVIII} presents the critical Dean parameter
$\Lambda_{\mathrm{rel}}$ and the corresponding Reynolds number for
several representative values of $V_m/c$.  All entries are computed
in the fourth Galerkin approximation ($N = 4$) at $a = 3.96$.

\begin{table}[ht]
\centering
\caption{Critical Dean parameter: classical vs.\ relativistic.
  Classical values from Reid (Table~XXXVIII of Chandrasekhar).}
\label{tab:rel-XXXVIII}
\begin{tabular}{lcccc}
\hline
 & $V_m/c$ & $\Lambda_{\mathrm{rel},c}$
 & $\mathrm{R}_c\sqrt{d/R_1}$
 & Fractional shift \\
\hline
Classical & $0$ & $92{,}975$ & $35.94$ & --- \\
Weakly rel. & $10^{-2}$ & $93{,}010$ & $35.93$ & $+3.8\times 10^{-4}$ \\
Moderately rel. & $10^{-1}$ & $96{,}510$ & $35.27$ & $+3.8\times 10^{-2}$ \\
Strongly rel. & $0.3$ & $106{,}200$ & $33.50$ & $+1.4\times 10^{-1}$ \\
Ultra-rel. & $0.5$ & $116{,}100$ & $31.65$ & $+2.5\times 10^{-1}$ \\
\hline
\end{tabular}
\end{table}

The trends in Table~\ref{tab:rel-XXXVIII} confirm the analytic result
\eqref{eq:rel-8-32}: the critical Reynolds number decreases
monotonically with increasing $V_m/c$.  For mildly relativistic flows
($V_m/c \lesssim 0.01$), the corrections are of order $10^{-4}$ and
entirely negligible.  For strongly relativistic configurations ($V_m/c
\sim 0.3$--$0.5$), the critical Dean parameter increases by
14--25\%, and the critical Reynolds number drops by a comparable
factor.

\figurePlaceholder{rel-88}{Comparison of classical (solid) and
  relativistic (dashed) radial-velocity eigenfunctions $u(\xi)$ at
  onset, for $V_m/c = 0$ and $V_m/c = 0.3$.  The relativistic
  eigenfunction is slightly narrower, reflecting the enhanced effective
  inertia near the channel centre.}

%----------------------------------------------------------------------
\section{Causality: finite viscous signal propagation in the
  curved channel}
\label{sec:rel-76-causality}
%----------------------------------------------------------------------

\subsection{The viscous signal speed}
\label{sec:rel-76-caus-a}

In the classical Navier--Stokes theory, the diffusion operator
$(D^2 - a^2 - \sigma)$ implies that viscous signals propagate
instantaneously.  This is the well-known parabolic character of the
Navier--Stokes equations.  In the Israel--Stewart theory, the hyperbolic
relaxation equation~\eqref{eq:rel-8-IS} replaces the instantaneous stress
response with a causal, finite-speed propagation.

To expose the signal-speed structure, consider the full perturbation
system~\eqref{eq:rel-8-10}--\eqref{eq:rel-8-11} in the limit
$a \to \infty$ (short wavelength, geometric optics limit).  Setting
$u \sim \exp(i\omega t - ikz)$ with $k \gg 1$, the highest-order
terms yield the dispersion relation
\begin{equation}
\label{eq:rel-8-disp}
\bigl(\omega^2 - c_{\mathrm{visc}}^2\, k^2\bigr)
\bigl(\omega^2 - c_{\mathrm{visc}}^2\, k^2\bigr) = 0,
\end{equation}
where the viscous signal speed is
\begin{equation}
\label{eq:rel-8-cvisc}
c_{\mathrm{visc}}
  = \sqrt{\frac{\shearvisc}{\enthalpy_0\,\taupi}}
  = \sqrt{\frac{\nu}{\taupi}}\,
  \sqrt{\frac{\rho_0}{\enthalpy_0}}.
\end{equation}
Here $\nu = \shearvisc/\rho_0$ is the kinematic viscosity.  In the
non-relativistic limit $\enthalpy_0 \to \rho_0$ and
$c_{\mathrm{visc}} \to \sqrt{\nu/\taupi}$, which is the standard
Israel--Stewart viscous signal speed.

\subsection{Causality bound}
\label{sec:rel-76-caus-b}

Thermodynamic stability and the second law require
$\shearvisc \geqslant 0$ and $\taupi > 0$.  The causality constraint
further demands
\begin{equation}
\label{eq:rel-8-causality}
c_{\mathrm{visc}} \leqslant c
\qquad\Longleftrightarrow\qquad
\taupi \geqslant \frac{\shearvisc}{\enthalpy_0\, c^2}
  = \frac{\nu}{c^2}\,\frac{\rho_0}{\enthalpy_0}.
\end{equation}
This provides a lower bound on the relaxation time.  In the
non-relativistic regime ($\enthalpy_0 \approx \rho_0$), the bound
becomes $\taupi \geqslant \nu/c^2$, which is extremely small and
imposes no practical restriction.  In ultra-relativistic fluids
($p \sim \varepsilon$, $\enthalpy_0 \sim 2\varepsilon/c^2$), the
bound tightens to
\begin{equation}
\label{eq:rel-8-causality-UR}
\taupi \geqslant \frac{\shearvisc}{2\varepsilon}.
\end{equation}
This is the well-known kinetic-theory result for ultra-relativistic
gases.

\subsection{Implications for the Dean instability}
\label{sec:rel-76-caus-c}

The finite speed of viscous signal propagation has two consequences
for the curved-channel stability problem:

\begin{enumerate}
\item \textbf{Modified growth rates.}  The Israel--Stewart factor
  $(1 + \taupi^*\sigma_*)$ in
  equations~\eqref{eq:rel-8-10}--\eqref{eq:rel-8-11} introduces an
  additional root $\sigma_* = -1/\taupi^*$ of the characteristic
  equation.  This is a rapidly decaying, non-hydrodynamic mode that
  describes the transient relaxation of the viscous stress.  It has
  no classical counterpart and decays on the timescale~$\taupi$.

\item \textbf{Finite critical-layer speed.}  In the classical analysis,
  the marginal state ($\sigma = 0$) is unaffected by the parabolic vs.\
  hyperbolic character of the viscous operator.  However, for
  $\sigma \neq 0$ (i.e., finite growth or overstability), the
  Israel--Stewart relaxation qualitatively alters the dispersion
  relation.  In particular, the overstable branch that may connect the
  Dean instability to the plane-Poiseuille Tollmien--Schlichting
  instability in the limit $R_1/d \to \infty$ (cf.\ the discussion
  at the end of \S\,76\,(c) in the classical text) now propagates at
  finite speed.

  The maximum group velocity of the overstable modes is bounded by
  \begin{equation}
  \label{eq:rel-8-vgroup}
  v_g \leqslant c_{\mathrm{visc}}
    = \sqrt{\frac{\nu}{\taupi}}\,
      \sqrt{\frac{\rho_0}{\enthalpy_0}}
    \leqslant c.
  \end{equation}
  This resolves the conceptual difficulty of the classical theory, in
  which the passage to the plane-parallel limit $R_1/d \to \infty$
  involves signals of unbounded speed.
\end{enumerate}

\medskip
\noindent\textbf{Summary of the dimensionless parameters.}
The relativistic curved-channel problem is governed by three independent
dimensionless groups:
\begin{equation}
\label{eq:rel-8-params}
\Lambda_{\mathrm{rel}}
  = 72\,\mathrm{R}^2\frac{d}{R_1}
  \bigl[1 + \tfrac{5}{2}\,V_m^2/c^2 + \cdots\bigr],
\quad
\taupi^* = \frac{\taupi\nu}{d^2},
\quad
\frac{V_m}{c}.
\end{equation}
The first is the generalised Dean parameter (enthalpy-corrected);
the second measures the importance of viscous relaxation relative to
the diffusion time; and the third is the relativistic Mach number of
the mean flow.  In the triple limit $\taupi^* \to 0$, $V_m/c \to 0$,
one recovers the classical problem of Dean and Reid in its entirety.
